% =====================================================================
% AI Conquer 2026 — Report (hướng ĐI LÀM / Application)
% Topic: Budget Planner — Module 01
% File khởi đầu cho team. Import cả thư mục này lên Overleaf, đặt file
% này làm Main File rồi điền nội dung. Giữ nguyên format chung.
% Cấu trúc bám mẫu Templates/Application_Template_Report.pdf
% =====================================================================
\documentclass[12pt]{ai_conquer2026}
\addbibresource{references.bib}

\title{%
  \begin{center}
        \Large {AI VIET NAM -- AI Conquer 2026}\\[.5ex]
        \Huge \textbf{Budget Planner \\ (hướng Đi làm)}\\[.5ex]
  \end{center}
}

% Author --------------------------------------------------------------
\posttitle{
    \author{
        \begin{minipage}{\textwidth}
        \centering
        \vspace{-1.5em}
        Tên nhóm / Thành viên: Đức Trần, Hồ Mai Thảo Nguyên,\\
        Phạm Viết Trường, Nguyễn Thị Phương, Lê Quang Cảnh
        \end{minipage}
    }
}
\date{}

% Header --------------------------------------------------------------
\pagestyle{fancy}
\fancyhf{}
\lhead{\href{https://www.facebook.com/aivietnam.edu.vn}{\textcolor{aivnGreen}{\bfseries AI VIET NAM (AI Conquer 2026)}}}
\rhead{\href{https://aivietnam.edu.vn/}{\textcolor{aivnGreen}{\bfseries aivietnam.edu.vn}}}
\fancyfoot[C]{\thepage}
\renewcommand{\headrulewidth}{1.0pt}
\renewcommand{\headrule}{{\color{aivnGreen}\hrule width\headwidth height\headrulewidth \vskip-\headrulewidth}}

\begin{document}
\maketitle
\vspace{-5.5em}

% --- I -----------------------------------------------------------------
\section{Tổng quan báo cáo}
% Project giải quyết bài toán gì, vì sao cần AI, sản phẩm cuối là gì,
% hệ thống hoạt động tổng quát ra sao. (Tên project, mục tiêu, vấn đề
% thực tế, lý do dùng AI/ML, người dùng mục tiêu, dữ liệu, phương pháp
% chính, kết quả nổi bật của demo.)
TODO: Budget Planner — ứng dụng theo dõi & cảnh báo chi tiêu...

\newpage
\setcounter{tocdepth}{2}
\tableofcontents
\thispagestyle{fancy}
\newpage

% --- II ----------------------------------------------------------------
\section{Định nghĩa bài toán}
% Bối cảnh vấn đề · Mục tiêu · Đầu vào hệ thống · Đầu ra mong muốn ·
% Loại bài toán AI · Tiêu chí thành công.
TODO

% --- III ---------------------------------------------------------------
\section{Dữ liệu sử dụng}
% Nguồn dữ liệu giao dịch, schema (ngày, category, số tiền, thu/chi),
% kích thước, cách thu thập/sinh, chất lượng.
TODO

% --- IV ----------------------------------------------------------------
\section{Tiền xử lý và chuẩn bị dữ liệu}
% Làm sạch, chuẩn hoá ngày/số tiền, xử lý thiếu, phân loại category.
TODO

% --- V -----------------------------------------------------------------
\section{Phương pháp và mô hình sử dụng}
% Rule-based và/hoặc ML phân loại giao dịch, logic ngân sách & cảnh báo.
TODO

% --- VI ----------------------------------------------------------------
\section{Kiến trúc hệ thống}
% Sơ đồ pipeline: Input -> Validation -> Preprocessing -> Tổng hợp &
% ngân sách -> Báo cáo/biểu đồ. (Đặt hình vào Figures/.)
TODO

% --- VII ---------------------------------------------------------------
\section{Luồng huấn luyện và suy luận}
\subsection{Luồng huấn luyện}
% (Nếu có ML phân loại/dự báo.) Nếu chỉ rule-based, nêu rõ và giải thích.
TODO
\subsection{Luồng suy luận}
TODO

% --- VIII --------------------------------------------------------------
\section{Công cụ và công nghệ sử dụng}
% Python, pandas/numpy, matplotlib, (streamlit), pytest, Git...
TODO

% --- IX ----------------------------------------------------------------
\section{Cấu trúc mã nguồn và hướng dẫn chạy}
% Cây thư mục repo + lệnh cài đặt + lệnh chạy demo.
TODO

% --- X -----------------------------------------------------------------
\section{Thực nghiệm và đánh giá mô hình}
% Kết quả, biểu đồ, độ chính xác phân loại (nếu có), demo input->output.
TODO

% --- XI ----------------------------------------------------------------
\section{Phân tích lỗi và hạn chế}
TODO

% --- XII ---------------------------------------------------------------
\section{Kết luận và hướng phát triển}
% Tổng kết + hướng mở rộng (Plan B: DB, UI, AI auto-categorize, deploy).
TODO

% --- XIII --------------------------------------------------------------
\section{Tài liệu tham khảo}
\label{section:references}
\nocite{*}        % bỏ dòng này nếu chỉ muốn in tài liệu đã trích dẫn
\printbibliography

\begin{appendices}
\begin{enumerate}
    \item \textbf{Mã nguồn:} Repo project tại \href{https://github.com/}{đây}. % TODO: link repo
    \item \textbf{Dataset:} Dữ liệu giao dịch mẫu tại \href{https://aivietnam.edu.vn/}{đây}. % TODO
\end{enumerate}
\end{appendices}

\end{document}
